%% kapa: exam=1 round=28 year=2017 period=1 form=A subject=real_estate subject_ko=부동산학원론 seq=30 number=110
%% source: Q-Net 감정평가사 기출문제 (cst00309) artlSeq=5199603 file=extracted/1008/5199603/2017년 제28회 감정평가사 1차시험 문제지/1교시 A형.pdf
%% conversion: kordoc 4.24.51 -> tools/gen1.py (auto); review=proofread-sonnet; flags=none
\begin{question}{110}
부동산의 수요와 공급, 균형에 관한 설명으로 옳은 것은? (단, 다른 조건은 동일함)
\begin{choices}
  \item 부동산의 수요는 유효수요의 개념이 아니라, 단순히 부동산을 구입하고자 하는 의사만을 의미한다.
  \item 건축비의 하락 등 생산요소 가격의 하락은 주택공급곡선을 왼쪽으로 이동시킨다.
  \item 수요자의 소득이 변하여 수요곡선 자체가 이동하는 경우는 수요량의 변화에 해당한다.
  \item 인구의 증가로 부동산 수요가 증가하는 경우 균형가격은 상승하고, 균형량은 감소한다.
  \item 기술의 개발로 부동산 공급이 증가하는 경우 수요의 가격탄력성이 작을수록 균형가격의 하락폭은 커지고, 균형량의 증가폭은 작아진다.
\end{choices}
\end{question}
\answer{5}
